\chapter{SEN12MS-CR Dataset and Preprocessing}
\label{chap:ch07}

\section{The SEN12MS-CR Benchmark Dataset}
The SEN12MS-CR (Sentinel-1 and Sentinel-2 Multi-Spectral Dataset for Cloud Removal) dataset, introduced by Ebel et al.~\cite{ebel2020sen12mscr}, is the global benchmark for multimodal satellite image reconstruction.

The dataset spans all meteorological seasons and continents, covering diverse biomes including dense Amazonian rainforests, arid deserts, European agricultural farmlands, urban megacities, and Arctic tundras.

\subsection{Structure of a Multimodal Triplet}
Each observation consists of three spatially co-registered data arrays acquired over the exact same geographical bounding box:
\begin{enumerate}
    \item \textbf{Input Sentinel-2 Optical Image ($\mathbf{X}_{\text{cloudy}}$):} Acquired on a cloudy day, containing 4 spectral bands (\texttt{B04}, \texttt{B03}, \texttt{B02}, \texttt{B08}) contaminated by cumulus, stratus, cirrus, or haze.
    \item \textbf{Input Sentinel-1 SAR Image ($\mathbf{X}_{\text{SAR}}$):} Acquired close in time to the cloudy optical acquisition, containing VV and VH microwave backscatter channels.
    \item \textbf{Ground-Truth Reference Optical Image ($\mathbf{Y}^*_{\text{clear}}$):} Acquired over the same scene on a nearby clear-sky date.
\end{enumerate}

\begin{warningbox}{The Reality of Satellite Ground-Truth: Temporal Offset}
In satellite remote sensing, \textbf{it is physically impossible to acquire a cloud-free optical image and a cloudy optical image at the exact same second in time!} 
The clear reference target is taken from a different satellite overpass (typically a few days or weeks apart). As a consequence:
\begin{itemize}
    \item Seasonal crop harvesting or plowing may have occurred between the two dates.
    \item Soil moisture or river water levels may have changed.
    \item The reference image may occasionally contain faint residual cirrus haze or cloud shadows.
\end{itemize}
A research-level engineer must recognize that the ground-truth target contains natural temporal variations, meaning that penalizing tiny high-frequency pixel differences too harshly can cause the network to overfit to temporal noise!
\end{warningbox}

\section{Dataset Splitting: Eliminating Spatial Leakage}
Our dataset is pre-processed and cached in the directory \texttt{cache\_fall\_v3} across three partitions:
\begin{itemize}
    \item \textbf{Training Set (\texttt{train}):} 33,468 patches.
    \item \textbf{Validation Set (\texttt{val}):} 1,115 patches.
    \item \textbf{Test Set (\texttt{test}):} 151 patches.
\end{itemize}
Crucially, these splits are partitioned by \textbf{Region of Interest (ROI) scene directory} rather than by individual patches. If patches from the same $100\text{ km} \times 100\text{ km}$ tile were randomly mixed, the model could easily memorize road networks and geographic coastlines, creating the illusion of generalization. Scene-level splitting guarantees that the validation and test patches come from geographical territories that the model has never encountered during training!

\section{Exact Radiometric Normalization Formulas}
Raw sensor readings have vastly different physical units and dynamic ranges. Feeding unnormalized values into a neural network destabilizes activation variances and causes vanishing or exploding gradients.

\subsection{Optical Reflectance Normalization}
Raw Sentinel-2 Level-1C and Level-2A data are distributed as 16-bit integer Digital Numbers ($\text{DN} \in [0, 10000]$). In our cached dataset:
\begin{equation}
\mathbf{X}_{\text{opt}} = \text{clip}\left(\frac{\text{DN}}{10000.0}, 0.0, 1.0\right)
\end{equation}
A value of $0.0$ corresponds to $0\%$ surface reflectance (absolute absorption), while $1.0$ corresponds to $100\%$ surface reflectance. In practice, dense clouds often saturate near $0.8\text{--}1.0$, while typical clear vegetation has reflectances of $0.03\text{--}0.15$ in visible bands and $0.30\text{--}0.50$ in the NIR band.

\subsection{Sentinel-1 SAR Normalization: \texttt{paper\_exact}}
Synthetic Aperture Radar backscatter is measured in logarithmic \textbf{decibels (dB)}:
\begin{equation}
\sigma^0_{\text{dB}} = 10 \log_{10}(\sigma^0)
\end{equation}
Typical terrestrial SAR backscatter values range from $-35\text{ dB}$ (smooth calm water or tarmac) to $0\text{ dB}$ (corner reflectors in urban buildings). In \texttt{data/cache\_dataset.py}, our active training pipeline applies the \texttt{paper\_exact} normalization:
\begin{align}
\mathbf{X}_{\text{VV, norm}} &= \frac{\text{clip}(\sigma^0_{\text{VV}}, -25.0, 0.0) + 25.0}{25.0} \in [0, 1] \\
\mathbf{X}_{\text{VH, norm}} &= \frac{\text{clip}(\sigma^0_{\text{VH}}, -35.0, 0.0) + 35.0}{35.0} \in [0, 1]
\end{align}

\begin{mathbox}{SAR Normalization Walkthrough}
Suppose a patch contains two pixels: Pixel A over a calm lake and Pixel B over an urban skyscraper.
\begin{enumerate}
    \item \textbf{Pixel A (Calm Lake):}
    \[
    \sigma^0_{\text{VV}} = -28.4\text{ dB}, \quad \sigma^0_{\text{VH}} = -38.1\text{ dB}
    \]
    Clamping: $\sigma^0_{\text{VV}} \to -25.0$, $\sigma^0_{\text{VH}} \to -35.0$.
    Normalized values:
    \[
    \mathbf{X}_{\text{VV, norm}} = \frac{-25.0 + 25.0}{25.0} = \mathbf{0.0}, \quad \mathbf{X}_{\text{VH, norm}} = \frac{-35.0 + 35.0}{35.0} = \mathbf{0.0}
    \]
    \item \textbf{Pixel B (Urban Skyscraper):}
    \[
    \sigma^0_{\text{VV}} = -5.0\text{ dB}, \quad \sigma^0_{\text{VH}} = -14.0\text{ dB}
    \]
    Clamping: Values are already within $[-25, 0]$ and $[-35, 0]$.
    Normalized values:
    \[
    \mathbf{X}_{\text{VV, norm}} = \frac{-5.0 + 25.0}{25.0} = \frac{20.0}{25.0} = \mathbf{0.80}
    \]
    \[
    \mathbf{X}_{\text{VH, norm}} = \frac{-14.0 + 35.0}{35.0} = \frac{21.0}{35.0} = \mathbf{0.60}
    \]
\end{enumerate}
Both radar and optical channels are now bounded within $[0, 1]$, enabling stable multimodal feature fusion.
\end{mathbox}

\begin{codebox}{Code Implementation in \texttt{data/cache\_dataset.py}}
\begin{lstlisting}[language=Python]
# Lines 63-75 in data/cache_dataset.py
if self.sar_norm == 'paper_exact':
    # VV: [-25, 0] dB -> [0, 1]
    # VH: [-35, 0] dB -> [0, 1]
    s1[0] = np.clip(s1[0], -25.0, 0.0)
    s1[0] = (s1[0] + 25.0) / 25.0
    s1[1] = np.clip(s1[1], -35.0, 0.0)
    s1[1] = (s1[1] + 35.0) / 35.0
\end{lstlisting}
\end{codebox}

\section{Summary of Chapter 7}
\begin{enumerate}
    \item SEN12MS-CR provides spatially co-registered triplets of cloudy optical, SAR, and clear-sky reference imagery across global biomes.
    \item Due to orbital revisit schedules, clear-sky ground-truth imagery is acquired on a nearby date, introducing subtle temporal phenology variations.
    \item Strict scene-level dataset splitting prevents spatial data leakage between training, validation, and test sets.
    \item Optical surface reflectance is normalized to $[0, 1]$ by dividing digital numbers by 10,000.
    \item Sentinel-1 SAR dB backscatter is clipped to $[-25, 0]\text{ dB}$ (VV) and $[-35, 0]\text{ dB}$ (VH) and scaled to $[0, 1]$.
\end{enumerate}

\section{Practice Exercises}
\begin{enumerate}
    \item If a radar backscatter reading for VH is recorded as $-42.0\text{ dB}$, what is its normalized value under the \texttt{paper\_exact} scheme?
    \item Why are different clipping thresholds used for Sentinel-1 VV ($-25\text{ dB}$) and VH ($-35\text{ dB}$)?
    \item A raw Sentinel-2 pixel in band \texttt{B08} has a Digital Number of $3{,}450$. What is its surface reflectance $\rho$?
\end{enumerate}

\subsection*{Solutions}
\begin{enumerate}
    \item Clamping $-42.0\text{ dB}$ to $[-35.0, 0.0]$ yields $-35.0\text{ dB}$. Normalized: $\frac{-35.0 + 35.0}{35.0} = 0.0$.
    \item VH is a cross-polarized signal (vertical transmit, horizontal receive). Due to physical polarization depolarization losses in microwave backscattering, VH return signals are naturally $6\text{--}10\text{ dB}$ weaker than co-polarized VV returns. Lowering the minimum threshold to $-35\text{ dB}$ preserves subtle vegetation volume scattering that would otherwise be clipped out.
    \item $\rho = \frac{3450}{10000.0} = 0.345$ (characteristic of moderate green vegetation canopy reflectance in NIR).
\end{enumerate}
